\chapter*{Abstract}

Retrieval-Augmented Generation (RAG) has emerged as a foundational paradigm for
grounding Large Language Model (LLM) responses in external, verifiable knowledge
sources. While considerable progress has been made in single-turn RAG settings, the
multi-turn conversational scenario — where a system must generate accurate, faithful
responses across an extended dialogue — remains significantly more challenging and
comparatively understudied. In multi-turn interactions, retrieval quality degrades across
turns due to non-standalone queries, pronoun coreferences, and implicit topic carryover,
while generation quality is further compromised by error accumulation, unanswerable
questions, and the growing complexity of conversational context.
 
This thesis presents a unified, stability-oriented architecture for multi-turn RAG,
developed and evaluated on the MTRAG benchmark as part of the SemEval-2026 Task 8
shared task. The system addresses all three subtasks: passage retrieval (Task A),
reference-grounded response generation (Task B), and end-to-end RAG (Task C). The
central design principle is retrieval stability through query diversity: rather than relying
on heterogeneous retriever ensembles, we issue five complementary LLM-based query
reformulations — Minimal, Corpus-Specific, Hypothetical Document Embedding (HyDE),
Chain-of-Thought (CoT), and Anchor-Keyword — to a single corpus-aligned sparse
retriever (ELSER v1), and aggregate the resulting ranked lists via a hierarchical
Nested Reciprocal Rank Fusion (RRF) scheme.
 
For response generation, we adopt an agentic, multi-stage pipeline that decomposes
grounded generation into five sequential stages: answerability triage, evidence
span extraction, dual-candidate drafting, judge-based selection, and a final micro-adjustment pass. A 4-gram
extractiveness shaping function enforces a faithfulness-naturalness trade-off, penalizing
both hallucination-prone under-extractive responses and robotic over-extractive ones.
The end-to-end RAG pipeline further incorporates a multi-judge answerability gate with
confidence-weighted voting to handle the critical challenge of unanswerable turns.
 
Ablation studies on the development set indicate that: (i) query diversity over a single
corpus-aligned retriever outperforms heterogeneous retriever ensembling under the
benchmark's official relevance judgments; (ii) evidence extraction before generation
improves faithfulness; (iii) answerability calibration becomes the decisive factor in end-to-end performance
when the share of unanswerable requests rises, while top-rank retrieval errors remain a
substantial second source of failure. Our system ranked 1st in Task A (nDCG@5: 0.5776, +20.5\% over the
strongest baseline), 2nd in Task B (Harmonic Mean: 0.7698), and 11th in Task C
(Harmonic Mean: 0.5409) on the official SemEval-2026 leaderboard, among 38, 26, and 29
participating systems respectively.
 
\paragraph*{Keywords ---}
Retrieval-Augmented Generation, Multi-Turn Conversation, Query Rewriting,
Reciprocal Rank Fusion, Agentic Generation, Answerability Detection,
Large Language Models, Information Retrieval, SemEval-2026.